% TOP_PROBLEM: {"id": 2510, "title": "Kourovka Notebook Problem 21.1", "category_id": 17, "category": {"id": 17, "name": "group_theory", "display_name": "Group Theory", "description": "Problems about groups, group actions, representations, and related algebraic structures.", "slug": "group-theory", "order_index": 17, "created_at": "2026-07-31T15:26:25.671Z"}, "status": "open", "problem_number": "KOU-21.1", "set_id": 10}
% FIRST_POSTED: 2026-10-01
% ENTRY_KIND: statement_only
% UnsolvedMath ID 2510; source catalogue code KOU-21.1
% !TeX program = lualatex
% Definition and sources only; this file is not an LLM solution attempt.
\documentclass[11pt,a4paper]{article}
\usepackage[margin=25mm]{geometry}
\usepackage{fontspec}
\setmainfont{lmroman10-regular.otf}[BoldFont=lmroman10-bold.otf,
 ItalicFont=lmroman10-italic.otf,BoldItalicFont=lmroman10-bolditalic.otf]
\usepackage{amsmath,amssymb,mathtools}
\usepackage{unicode-math}
\setmathfont{latinmodern-math.otf}
\AtBeginDocument{\let\setminus\smallsetminus}
\usepackage{xurl,hyperref}
\hypersetup{colorlinks=true,urlcolor=blue}
\setcounter{secnumdepth}{0}
\setlength{\parindent}{0pt}
\setlength{\parskip}{5pt}
\setlength{\emergencystretch}{3em}
\begin{document}
\section*{Kourovka Notebook Problem 21.1}\label{2510}
{\small UnsolvedMath ID 2510 · KOU-21.1 · Group Theory · Kourovka Notebook - New Problems, Issue 21}
\subsection{Definitions and mathematical statement}
For a positive integer $n$, a finite group $K$ and an automorphism
$\varphi$ with $\varphi^n=\mathrm{id}$, use $x^{\varphi^j}=\varphi^j(x)$
and define the twisted solution set
\[
 X_{n,\varphi}(K)=\{x\in K:x\varphi(x)\cdots\varphi^{n-1}(x)=1\}.
\]
The factors are multiplied in the displayed order. Define $c_n$ as the
supremum of $|X_{n,\varphi}(K)|/|K|$ over all these pairs for which
$X_{n,\varphi}(K)\ne K$. Thus $c_n<1$ asks for one positive gap from
one, uniform over all finite groups and admissible automorphisms, after
excluding pairs where the equation holds identically.

Part (a) asks: for every $n>1$, does the hypothesis $c_{p^a}<1$ for
every prime power $p^a$ dividing $n$, with $a\ge1$, imply $c_n<1$?
The bounds for different prime powers need not be identical.

For part (b), put
\[
 H_n(G)=\langle\{x\in G:x^n\ne1\}\rangle
\]
for each finite group $G$, where brackets denote the subgroup generated
by the indicated elements; the empty set generates $1$.
Suppose there exists an integer $k_n\ge1$ depending only on $n$ such
that $[G:H_n(G)]\le k_n$ whenever $H_n(G)\ne1$.
Does this uniform index hypothesis imply $c_n<1$?

The index condition in (b) involves ordinary powers and arbitrary finite
$G$, whereas its conclusion controls twisted powers for arbitrary
$\varphi$. Both parts ask implications for a fixed integer parameter,
not bounds depending on the individual group.
\subsection{Short English statement}
Do prime-power density gaps, or a uniform generalized Hughes–Thompson index bound, force a uniform gap for twisted $n$th-power solutions?
\subsection{Sources}
\begin{itemize}
\item[S1] ULAM.AI and the UnsolvedMath Contributors, supplied problems.json, record 2510 (KOU-21.1), Kourovka Notebook - New Problems, Issue 21. Catalogue identity and source transcription; CC BY 4.0.
\url{https://huggingface.co/datasets/ulamai/UnsolvedMath}.
\item[S2] The Kourovka Notebook, 21st edition, new problems, Problem 21.1. Expanded formulation of the supplied catalogue statement.
\url{https://arxiv.org/pdf/1401.0300v47}.
\end{itemize}
\end{document}
